% LUT-native multimodal seizure detection on SeizeIT2 -- Phase 1 draft.
% Every number in this file must come from a file in results/ produced by a committed script.
% Result tables are \input from paper/tables/*.tex (scripts/make_tables.py). Hardware results
% are placeholders for Phase 2 and are marked \PHASETWO.
\documentclass[conference]{IEEEtran}
\usepackage{amsmath,amssymb,graphicx,booktabs,xcolor,cite}
\newcommand{\PHASETWO}[1]{\textcolor{red}{[Phase 2: #1]}}
\newcommand{\TODO}[1]{\textcolor{blue}{[TODO: #1]}}
\newcommand{\RESULT}[1]{\textcolor{teal}{[result: #1]}}

\begin{document}
\title{Multiplier-Free Seizure Detection with LUT-Native Neural Networks on Behind-the-Ear EEG, ECG, EMG and Motion}
\author{\IEEEauthorblockN{\TODO{authors}}\IEEEauthorblockA{National Institute of Technology Warangal, India}}
\maketitle

\begin{abstract}
Wearable seizure detection must run for days on a small battery, which favours classifiers
without multipliers. We study LUT-native neural networks, whose neurons are FPGA lookup tables,
for patient-independent detection of focal seizures from two-channel behind-the-ear EEG, single-lead
ECG, deltoid EMG and chest motion in the SeizeIT2 dataset (125 patients, 11,626 hours).
All features are computed by an integer, causal, multiplier-free front end, encoded by thermometer
comparators, and classified by a differentiable weightless neural network (DWN) whose prediction is a
popcount difference. On the official split, \RESULT{headline detection result vs.\ the official
baselines, SzCORE event metrics, mean $\pm$ std over 5 seeds}, at an estimated \RESULT{LUT count}
LUTs and zero multipliers in the classifier. \TODO{one sentence on the strongest new method (N1--N6)}.
Hardware measurements are left to future work.
\end{abstract}

\section{Introduction}
Long-term monitoring with wearable sensors is the most practical way to document focal seizures
outside the epilepsy monitoring unit, but the devices have little energy for computation. Behind-the-ear
(bte) EEG is unobtrusive yet sees only two channels, so focal seizures that are clear on full-scalp
EEG can be faint or invisible on the wearable \cite{vandecasteele2020}. Adding ECG, EMG and motion
helps \cite{vandecasteele2021,bhagubai2023}, but most detectors that use them rely on multiply-heavy
models.

LUT-native neural networks replace weighted sums by lookup tables: each neuron is a Boolean function
of a few input bits that maps one-to-one onto an FPGA LUT. Differentiable weightless neural networks
(DWNs) \cite{bacellar2024} and differentiable logic-gate networks \cite{petersen2022} train such
networks with gradient descent and reach strong accuracy at a few thousand LUTs on tabular and image
benchmarks; tree ensembles can be mapped to LUTs in the same spirit \cite{khataei2025}. They have
not been studied for seizure detection, and clinical biosignals bring problems that benchmark datasets
do not: extreme class imbalance (0.12\,\% seizure time), patient-to-patient variability, missing
modalities, labels taken from a different sensor, and event-level clinical scoring.

\textbf{Contributions.}
(i)~The first differentiable LUT-native seizure detectors and the first benchmark of the LUT-native
family (DWN, logic-gate networks, TreeLUT) against standard machine learning on the official SeizeIT2
split, with SzCORE event-based scoring \cite{dan2024}.
(ii)~A fully integer, causal, multiplier-free multimodal feature front end with a bit-exact golden model,
a fixed-point specification and a reference Verilog implementation of the classifier that matches the
golden model on all test vectors.
\TODO{(iii)--(iv): the two or three Tier-2 methods that win on the inner folds (Section~\ref{sec:methods}),
stated only if they pass the paired test.}

\section{Related Work}
\textit{Wearable bte-EEG seizure detection.} Vandecasteele et al.\ showed that two bte channels
capture many focal seizures \cite{vandecasteele2020}; the SeizeIT2 dataset paper reports SVM and
ChronoNet baselines on its official split: 71.1\,\% sensitivity at 11.0 false alarms per hour and
84.2\,\% at 100.5 per hour \cite{bhagubai2025}. SeizFt won the 2023 wearable-EEG challenge with
CatBoost on engineered features \cite{alhussaini2023}.
\textit{Multimodal detection.} ECG adds value to limited-channel EEG \cite{vandecasteele2021,bhagubai2023};
EMG and motion target motor seizures \cite{beniczky2018,leijten2018}.
\textit{Hardware detectors.} Busia et al.\ run a spiking network on a low-power FPGA for patient-specific
detection on scalp EEG (CHB-MIT) \cite{busia2025}; BrainFuseNet fuses EEG, PPG and accelerometry for
edge deployment \cite{ingolfsson2024}.
\textit{LUT-native networks on biosignals.} Logic-gate and LUT networks have been applied to inter-patient
ECG arrhythmia classification on an Artix-7 FPGA \cite{mommen2026} and logic-gate networks to EEG
dementia and emotion classification on a CPU \cite{dharia2026}; recurrent logic-gate networks have been
proposed for sequence modelling \cite{buhrer2025}; the cost of thermometer encoding in DWN accelerators
has been measured \cite{mecik2025}. None addresses seizures, multimodal input, missing modalities or
event-level clinical metrics.

\section{Dataset and Protocol}
\textit{Data.} SeizeIT2 (OpenNeuro ds005873 v1.1.0, CC0) \cite{bhagubai2025} contains 125 patients with
focal epilepsy from five European epilepsy monitoring units recorded with the Byteflies Sensor Dot:
2-channel bte-EEG (all patients), single-lead chest ECG, left-deltoid EMG and chest-module accelerometer
and gyroscope. Our audit of the release (Table~\ref{tab:audit}) found 2,850 recordings, 11,626 hours of
EEG, 883 annotated seizures (702 in the 96 training patients, 181 in the 29 test patients), EEG/ECG/EMG
at 256\,Hz and motion at 25\,Hz, ECG/EMG missing in 3.1\,\% and motion in 3.9\,\% of the EEG time.
Seizures are annotated from video and full-scalp EEG, so some are not visible on the wearable; the release
does not mark which, and it does not contain centre or age per patient.
\textit{Split.} We use the official split exactly: training sub-001 to sub-096, testing sub-097 to sub-125
(the official ``validation'' set; there is no separate test set). All design choices, hyperparameters,
thresholds and post-processing are tuned with 5-fold patient-wise cross-validation inside the training
patients; each final model is scored once per seed on the test patients.
\textit{Scoring.} Primary: SzCORE event-based scoring \cite{dan2024} computed with the \texttt{timescoring}
library (any overlap, 30\,s pre- and 60\,s post-ictal tolerance, events closer than 90\,s merged, events
longer than 5\,min split), averaged over patients, with false positives per 24\,h. For comparison with the
dataset paper we also report its own metrics (any-overlap sensitivity, false alarms per hour, AUROC, AUPR),
computed with a port of the authors' evaluation code; the two metric families are never mixed.
\textit{Statistics.} Per-patient paired Wilcoxon tests against each baseline, Friedman test across models,
95\,\% patient-bootstrap intervals, mean $\pm$ standard deviation over 5 seeds.

\section{Method}\label{sec:methods}
\subsection{Integer multimodal front end}
Samples arrive as int16 at a fixed scale (0.25\,$\mu$V per count for biopotentials). Every filter is a
second-order section with coefficients rounded to at most five signed powers of two, implemented in direct
form~I with second-order error feedback: without it, the rounding error of the 0.5\,Hz high-pass is
amplified about 1,700 times at DC and exceeds the EEG itself. Features are computed every second over
2\,s windows (512 samples, so every mean is a shift): per EEG channel the mean absolute value, line length,
a log-domain Hjorth-mobility proxy, zero crossings, local extrema and four band magnitudes; symmetric
cross-channel combinations (mean, maximum, absolute difference) make the model insensitive to the
patient-specific montage; ECG heart rate and RR variability from a multiplier-free R-peak detector, in the
log domain so no division is needed; EMG magnitude, line length, zero crossings and a high/low band
ratio; chest accelerometer and gyroscope magnitude, jerk, activity count and 1--5\,Hz band magnitude;
per-modality ``present'' bits; and, for each magnitude feature, its difference from a causal 64\,s
moving average (a shift). Magnitudes are log-compressed by a leading-one detector with three mantissa bits.
The artefact flag replaces the official RMS rule by the mean absolute value against the same limits scaled
by $\sqrt{2/\pi}$, which agrees with the float rule on \RESULT{agreement}. The full front end uses no
multipliers (\texttt{spec/fixed\_point.md}).
\subsection{Thermometer encoding and DWN}
Each feature is compared against $b$ thresholds placed at training quantiles; each distinct threshold is
one comparator. A DWN layer of LUT6 neurons reads these bits through a learnable mapping and is trained
with extended finite differences \cite{bacellar2024}; a GroupSum over two classes gives the seizure score
as the popcount difference of the seizure and background neuron groups. A decision is made every second
when the score reaches a threshold and the artefact gate is clear; an alarm is raised when at least $k$ of
the last 10 decisions are positive.
\subsection{New methods}
\TODO{Describe only the methods kept after the inner-fold decision (results/decisions.md): N1 stateful LUT
network, N2 comparator-budgeted encoding, N3 label-free per-patient threshold adaptation, N4 distillation,
N5 seizure-level multiple-instance learning and noise-robust losses, N6 reliability-gated fusion,
N7 interpretable alarms.}

\section{Results}
\RESULT{Table: detection results on the 29 test patients (dataset-paper metrics and SzCORE).}
\RESULT{Table: LUT-native family benchmark with estimated LUTs, comparators, multipliers.}
\RESULT{Table: new-method ablations with differences, 95\,\% intervals and p-values.}
\RESULT{Table: modalities and missing-modality robustness.}
\RESULT{Table: results by seizure type (by centre and age group not possible with the public release).}
\RESULT{Figure: event sensitivity vs.\ false alarms per 24\,h; Pareto plot of event-F1 vs.\ estimated LUTs.}

\section{Hardware}
\PHASETWO{post-synthesis LUT/FF/BRAM counts on KV260 and Artix-7, power with switching activity,
bit-exact replay of the test recordings, comparison with the INT8 MLP through the same flow.}
Phase~1 reports only estimated costs computed from the exported models (Section~VI tables, labelled
``estimated'').

\section{Limitations}
Results come from retrospective replay of recordings, not live patients. Annotations come from full-scalp
EEG and video, so some seizures are invisible on the wearable channels, and the release does not identify
them. One dataset is used. Centre and age are not in the public release, so the folds cannot be stratified
by centre and no centre or adult/child breakdown is reported. Hardware costs are software estimates until
Phase~2.

\section*{Reproducibility and Ethics}
Code, configurations, export packages and the Verilog generator are released with the paper. Dataset:
SeizeIT2, DOI 10.18112/openneuro.ds005873.v1.1.0. The original study was approved by the UZ Leuven ethics
committee (S63631; data sharing S67350).

\bibliographystyle{IEEEtran}
\bibliography{refs}
\end{document}
